\section{An independent rank-four census}\label{sec:census}
We give a complete integer enumeration, with no input from an external
catalogue. In this appendix all parameters lie in $\{0,\ldots,M\}$.
The unit coefficients ensure that the multiplicity is the maximum of one
and the parameters. Commutativity follows from Section 3 of the accompanying completeness.md; duality has just two possible types.

\subsection{Self-dual parameters and elimination}
Frobenius reciprocity gives the following general matrices:
\[
N_x=\begin{pmatrix}0&1&0&0\\1&a&b&c\\0&b&d&e\\0&c&e&f\end{pmatrix},\quad
N_y=\begin{pmatrix}0&0&1&0\\0&b&d&e\\1&d&g&h\\0&e&h&i\end{pmatrix},\quad
N_z=\begin{pmatrix}0&0&0&1\\0&c&e&f\\0&e&h&i\\1&f&i&j\end{pmatrix}.
\]
Their pairwise commutativity is equivalent to the six equations
\begin{align}
 ad-b^2+bg+ch-d^2-e^2+1&=0,\label{eq:census1}\\
 ae-bc+bh+ci-de-ef&=0,\label{eq:census2}\\
 be-cd+dh-eg+ei-fh&=0,\label{eq:census3}\\
 af+bi-c^2+cj-e^2-f^2+1&=0,\label{eq:census4}\\
 bf-ce+di-eh+ej-fi&=0,\label{eq:census5}\\
 df-e^2+gi-h^2+hj-i^2+1&=0.\label{eq:census6}
\end{align}
These are necessary and sufficient for associativity. Indeed the first
row of $N_uN_v$ records $uv$. Commutativity implies that every other row
of $N_uN_v-\sum_wN_{uv}^wN_w$ is obtained by multiplying its zero first
row by a regular matrix, so the difference vanishes.

The six mixed coordinates $b,c,d,f,h,i$ form one orbit under
permutations of $x,y,z$. Every based-ring orbit therefore has a
representative with
\[
 b=\min\{b,c,d,f,h,i\}.
\]
For $b>0$, enumerate $b,c,d,e,f$, with $c,d,f\geq b$, and put
\begin{align*}
 H_0&=e(d+f)+bc,&G_0&=b^2+d^2+e^2-1,\\
 J_0&=c^2+e^2+f^2-1,&
 A&=e(b^2-c^2)+bc(f-d),\\
 C&=e(bf-ce),&
 B_0&=b^3e-b^2cd+b(d-f)H_0-ebG_0+ecH_0.
\end{align*}
The first, second and fourth equations give
\[
 h=\frac{H_0-ae-ci}{b},\qquad
 g=\frac{G_0-ad-ch}{b},\qquad
 j=\frac{J_0-af-bi}{c}.
\]
After substitution, the third equation multiplied by $b^2$ is
\begin{equation}\label{eq:census-linear}
 Ca+Ai+B_0=0.
\end{equation}
If $A\ne0$, let $q=\gcd(A,C)$. There is no solution unless $q\mid B_0$.
Otherwise $a$ runs through the unique residue class solving
$(C/q)a\equiv-B_0/q\pmod{|A/q|}$, and $i$ is determined.
A modulus of one means every bounded $a$. If $A=0\ne C$, determine
$a$ and enumerate $i$. If $A=C=0$, require $B_0=0$ and retain all bounded
$a,i$. Impose $0\leq a,g,j\leq M$ and $b\leq h,i\leq M$,
exact divisibility, and all six original equations.

If $b=0<c$, determine
\[
 h=\frac{d^2+e^2-ad-1}{c},\quad
 i=\frac{e(d+f-a)}c,\quad
 j=\frac{c^2+e^2+f^2-af-1}c.
\]
For $e>0$, integrality of $i$ gives
$a\equiv d+f\pmod{c/\gcd(c,e)}$. Enumerate $c,d,e,f$ and this
progression for $a$; the bounds on $h,i,j$ give further linear interval
restrictions. The equation $eg=h(d-f)+ei-cd$ determines $g$.
For $e=0$, nonnegativity forces $d,h>0$ and
\[
 \gcd(c,d)=1,\qquad a\equiv d-d^{-1}\pmod c,\qquad
 f=d-\frac{cd}{h}.
\]
Indeed $ch=d^2-ad-1$ implies the coprimality, and $h(d-f)=cd$
gives the last identity. Enumerate $c,d,a$, determine $h,f,j$, and
retain every bounded $g$. The congruence modulo one is unrestricted.

Finally, $b=c=0$ gives
\[
 ad-d^2-e^2+1=0,\quad e(a-d-f)=0,\quad af-e^2-f^2+1=0.
\]
For $e>0$ these are $a=d+f$ and $df=e^2-1$. Enumerate the bounded
factor pairs, including $e=1,d=0$ with every bounded $f$, then enumerate
$h,i$ and determine
\[
 eg=h(d-f)+ei,\qquad ej=eh+i(f-d).
\]
For $e=0$, the preliminary equations force $a=0,d=f=1$. The remaining
equation is $gi+hj=h^2+i^2-2$. If $h>0$, enumerate $h,i,g$ and solve
for $j$; if $h=0<i$, solve for $g$ and retain every bounded $j$.
The case $h=i=0$ is impossible. In every branch, all six equations and
all coefficient bounds are tested again.

\subsection{One dual pair}
If $X^*=X$ and $Y^*=Z$, the general matrices are
\[
N_X=\begin{pmatrix}0&1&0&0\\1&a&b&b\\0&b&c&d\\0&b&d&c\end{pmatrix},\quad
N_Y=\begin{pmatrix}0&0&1&0\\0&b&c&d\\0&d&e&f\\1&c&e&e\end{pmatrix},\quad
N_Z=N_Y^{\mathsf T}.
\]
Associativity is equivalent to
\begin{align*}
 ac-b^2+2be-c^2-d^2+1&=0,&ad-b^2+be+bf-2cd&=0,\\
 b(c-d)+d(e-f)&=0,&-c^2+d^2-e^2+f^2-1&=0.
\end{align*}
Enumerate $c,d,e$. The last equation determines the unique nonnegative
candidate $f=\sqrt{1+c^2-d^2+e^2}$; retain it only if integral and
bounded. If $c\ne d$, determine $b=-d(e-f)/(c-d)$. If $c=d$, require
$d(e-f)=0$ and retain every bounded $b$. Determine $a$ from the first
equation if $c>0$, or from the second if $d>0$; if $c=d=0$, retain
all bounded $a$. Finally test all four equations and all bounds.
The square-root implementation corrects its initial estimate by integer
comparisons and checks its square exactly.

\begin{proposition}[Completeness and cost]\label{prop:census}
The preceding algorithm enumerates all rank-four fusion rings of
multiplicity at most $M$, up to based-ring isomorphism. The parameter
search uses $O((M+1)^6)$ integer operations before canonicalization.
\end{proposition}
\begin{proof}
The two presentations exhaust duality and reciprocity. Each branch either
solves a necessary linear equation by exact division or exhausts every
allowed value when its coefficient vanishes. Final equation checks are
necessary and sufficient for associativity. Thus no ring is omitted or
spuriously admitted. In the self-dual search choose the lexicographically
least parameter tuple among its permutations satisfying the same
minimal-$b$ condition. Comparing only those images is essential: the
unconstrained least tuple need not satisfy that condition. Each orbit has
exactly one retained tuple. In the other duality type the only possible
nontrivial based automorphism interchanges $Y,Z$, and fixes all six
parameters. Thus every retained non-self-dual tuple is already one class.
The full-tensor interface additionally canonicalizes under all six
permutations and checks distinctness independently.

There are five outer parameters in the positive-$b$ branch. Solving
\eqref{eq:census-linear} takes $O(\log(M+1))$ arithmetic operations for
the greatest common divisor and at most $M+1$ candidates. The apparently
singular branch $A=C=B_0=0$ cannot occur there: if $e=0$, then
$A=0$ forces $f=d$ and $B_0=-b^2cd\ne0$; if $e>0$, the equations
$A=C=0$ give $cd=be$, $bf=ce$, and substitution gives $B_0=be\ne0$.
The zero-$b$ branches use at most five free coordinates, and the
non-self-dual search at most five. The stated $O((M+1)^6)$ arithmetic
bound follows. Bit costs grow with $\log(M+1)$; the implementation
restricts its integer range explicitly.
\end{proof}

\subsection{Current implementation and complete run}
The current implementation is \texttt{code/rank4\_census.cpp}; the
reproduction entry point is \texttt{code/run\_census.py}. The full
run through multiplicity $128$ produces $889530$ based-isomorphism
classes. Every output tensor is independently checked for the full
fusion-ring axioms and canonicalized over all six unit-fixing
permutations. The exact counts and execution reports are in the
accompanying release. Exhaustiveness rests on the argument above,
not on comparison with an existing list.
